\documentclass[../../../include/open-logic-section]{subfiles}

\begin{document}

\olfileid[tr]{sth}{choice}{countablechoice}
\olsection{Sayılabilir Seçim}

Seçim/İyi Sıralama hiç kullanılmadan şunu ispatlamak kolaydır:

\begin{lem}[$\Zminus$ içinde]
Her sonlu kümenin bir seçim fonksiyonu vardır.
\end{lem}

\begin{proof}
$a=\{b_1,\ldots,b_n\}$ olsun. Kolaylık için her
$b_i\neq\emptyset$ varsayın. Öyleyse
$c_1\in b_1,\ldots,c_n\in b_n$ olan $c_1,\ldots,c_n$ nesneleri vardır.
Şimdi \olref[z][pairs]{prop:pairsconsequences} gereği
$\{\langle b_1,c_1\rangle,\ldots,\langle b_n,c_n\rangle\}$ kümesi
vardır; bu da $a$ için bir seçim fonksiyonudur.
\end{proof}

Fakat sonsuz kümeleri ele alır almaz işler bulanıklaşır. Örneğin
yukarıdakinin şu ``en küçük'' genişletmesini ele alın:

\begin{defish}
\emph{Sayılabilir Seçim.} Her \emph{sayılabilir} kümenin bir seçim
fonksiyonu vardır.
\end{defish}

Bu, Seçim'in özel bir durumudur. Üstelik ilkenin kullanıldığı açıkça fark
edilmeden oldukça sık başvurulduğu ortaya çıkar. İşte iki güzel
örnek.\footnote{\citet[\S9.4]{Potter2004} ve Luca Incurvati'ye borçluyuz.}

\begin{ex}
Doğal bir düşünce şudur: her $A$ kümesi için ya $\cardle{\omega}{A}$ ya
da bir $n\in\omega$ için $\cardeq{A}{n}$ olur. Bu, her kümenin ya sonlu
ya da sonsuz olduğu yönündeki sezgisel düşünceyi ifade etmenin bir yoludur.
Cantor ve başka birçok matematikçi bu iddiayı ispatlamadan ileri sürdü.
Biz temkinli davranarak bunu
\olref[cardinals][classing]{generalinfinitycharacter} içinde ispatladık.
Fakat o ispatta her $A$ kümesinin iyi sıralanabileceğini ve dolayısıyla
$\card{A}$'nın varlığının güvence altında olduğunu varsaydığımız için
$\ZFC$ içinde çalışıyorduk. Yani Seçim'i açıkça varsaydık.

Gerçekten \citet{Dedekind1888} bu iddia için kendi ispatını şöyle sundu:

\begin{thm}[$\Zminus + \text{Sayılabilir Seçim}$ içinde]
Her $A$ için ya $\cardle{\omega}{A}$ ya da bir $n\in\omega$ için
$\cardeq{A}{n}$ olur.
\end{thm}

\begin{proof}
Her $n\in\omega$ için $\cardneq{A}{n}$ varsayın. O zaman özellikle her
$n<\omega$ için tam $\cardexpo{2}{n}$ elemanlı bir
$A_n\subseteq A$ altkümesi vardır. Bu $A_0,A_1,A_2,\ldots$ dizisini
kullanarak her $n$ için şunu tanımlayın:
\[
	B_n = A_n \setminus \bigcup_{i < n} A_i.
\]
Şuna dikkat edin:
\begin{align*}
	\card{\bigcup_{i < n}A_i}
	&\leq \card{A_0} + \card{A_1} + \ldots + \card{A_{n-1}}\\
	&=1 + 2 + \ldots + 2^{n-1}\\
	& = 2^n - 1\\
	& < 2^n = \card{A_n}
\end{align*}
Dolayısıyla her $B_n$'nin en az bir $c_n$ elemanı vardır. Üstelik
$B_n$'ler ikişer ikişer ayrıktır; böylece $c_n=c_m$ ise $n=m$ olur.
Fakat her $c_n\in A$ olduğundan $f(n)=c_n$ fonksiyonu
$\omega\to A$ bir iğnelemedir.
\end{proof}
\noindent
Dedekind, Sayılabilir Seçim'i kullandığını belirtmedi. Peki
\emph{siz} kullanıldığı yeri gördünüz mü? Yeniden bakın. (Gerçekten:
\emph{yeniden bakın}.)

İspat Sayılabilir Seçim'i iki kez kullandı. Bir kez
$A_0$, $A_1$, $A_2$, \dots\@ kümeler dizisini elde etmek için kullandık. Sonra her
$B_n$'den $c_n$ elemanını seçmek için yeniden kullandık. Üstelik Seçim'in
bu kullanımı ortadan kaldırılamaz. \citet[p.~138]{Cohen1966}, Seçim'in
hiçbir sürümü olmadığında sonucun başarısız olduğunu ispatladı. Yani
$\omega$ ile \emph{karşılaştırılamayan} kümelerin bulunması $\ZF$ ile
tutarlıdır.
\end{ex}

\begin{ex}
\citeyear{Cantor1878} yılında Cantor, sayılabilir kümelerin sayılabilir
birleşiminin sayılabilir olduğunu belirtti. İspat sunmadı; belki de ispatın
açık olduğunu düşünüyordu. Biz temkinli davranarak bu sonucun daha genel
bir sürümünü \olref[card-arithmetic][simp]{kappaunionkappasize} içinde
ispatladık. Fakat ispatımız açıkça Seçim'i varsaydı. Daha az genel sonucun
ispatı bile Sayılabilir Seçim'i gerektirir.

\begin{thm}[$\Zminus + \text{Sayılabilir Seçim}$ içinde]
Her $n\in\omega$ için $A_n$ sayılabilir ise
$\bigcup_{n<\omega}A_n$ sayılabilirdir.
\end{thm}

\begin{proof}
Genelliği kaybetmeden her $A_n\neq\emptyset$ varsayın. Öyleyse her
$n\in\omega$ için bir $f_n\colon\omega\to A_n$ !!a{surjection} vardır.
$f(m,n)=f_n(m)$ ile
$f\colon\omega\times\omega\to\bigcup_{n<\omega}A_n$ tanımlayın.
$\omega\times\omega$ sayılabilir
(\olref[sfr][siz][zigzag]{natsquaredenumerable}) ve $f$ bir
!!a{surjection} olduğundan sonuç çıkar.
\end{proof}
\noindent
Sayılabilir Seçim'in kullanıldığı yeri gördünüz mü? $f_0$, $f_1$, $f_2$, \dots
fonksiyonlar dizisini seçmek için kullanılır.\footnote{Seçim'in benzer bir
kullanımı, \olref[card-arithmetic][simp]{kappaunionkappasize} içinde
``Her $\beta\in\cardfont{a}$ için bir $f_\beta$ !!a{injection}
sabitleyin'' yönergesini verdiğimizde ortaya çıktı.} Ve yine, herhangi bir
Seçim ilkesi olmadığında sonuç başarısız olur. Özellikle
\citet{FefermanLevy1963}, sayılabilir kümelerin sayılabilir bir birleşiminin
kardinalitesinin $\beth_1$ olmasının $\ZF$ ile tutarlı olduğunu ispatladı.
Fakat Russell aynı noktayı çok daha eğlenceli biçimde anlatır:
\begin{quote}
  Bu, her çizme çifti aldığında bir çift de çorap alan, başka hiçbir zaman
  çorap almayan ve her ikisini almaya öylesine tutkuyla bağlı olan bir
  milyonerle örneklenir ki sonunda $\aleph_0$ çift çizmesi ve $\aleph_0$
  çift çorabı olur\dots\@ Çizmeler arasında sağ ve solu ayırt edebiliriz;
  dolayısıyla her çiftten birini seçebiliriz: bütün sağ çizmeleri ya da
  bütün sol çizmeleri seçebiliriz. Fakat çoraplarda buna benzer hiçbir
  seçim ilkesi kendini göstermez ve çarpımsal aksiyomu [yani aslında
  Seçim'i] varsaymadıkça her çiftten bir çoraptan oluşan herhangi bir
  sınıfın bulunduğundan emin olamayız.
  \citep[p.~126]{Russell1919}
\end{quote}
Kısacası şunu ispatlamak için Seçim'in bir biçimi gerekir: sayılabilir
çoklukta çorap çiftiniz varsa (yalnızca) sayılabilir çoklukta çorabınız
vardır. Gerçekten Sayılabilir Seçim (ya da ona denk bir şey) olmadan
sayılabilir kümelerin sayılabilir bir birleşimi sayılabilir olmayabilir.
\end{ex}

Çıkarılacak ders, Sayılabilir Seçim'in kullanıcıları bunu pek fark etmeden
defalarca kullanılmış olmasıdır. Felsefi soru şudur: Sayılabilir Seçim'i
nasıl \emph{gerekçelendirebiliriz}?

Sezgisel bir gerekçelendirme girişimi bir süper göreve başvurabilir. İlk
seçimi bir dakikanın $\nicefrac{1}{2}$'sinde, ikinci seçimi
$\nicefrac{1}{4}$'ünde, \dots, $n$inci seçimi
$\nicefrac{1}{2^n}$'inde, \dots\@ yaptığımızı varsayın. O zaman $1$ dakika
içinde bir $\omega$-seçimler dizisi yapmış ve bir seçim fonksiyonu
tanımlamış oluruz.

Fakat böyle bir düşünce deneyi gerçekte bize ne söyleyebilir? Öncelikle bu,
``seçme'' düşüncesini oldukça kelimesi kelimesine almaya dayanır. Ayrıca
matematiği metafizik olanaklılığa bağlar görünmektedir.

Daha önemlisi: bize \emph{salt} Sayılabilir Seçim yerine \emph{koşulsuz}
Seçim için hiçbir gerekçe vermeyecektir. Çünkü \emph{her} kümenin bir
seçim fonksiyonu olması gerekiyorsa ``keyfî sıra sayısı uzunluğunda bir
süper görev'' gerçekleştirebilmemiz gerekir. Açıkçası bu düşünce
gülünçtür.

\end{document}
